\section{Research Design and Methodology}
\label{ch:methodology}

\subsection{Research strategy}

The project used an applied design-science strategy. Design science is appropriate when knowledge is generated through creating and evaluating an artefact that addresses an identified problem \parencite{hevner2004}. \textcite{peffers2007} describe a cycle of problem identification, objectives, design, demonstration, evaluation and communication. That cycle maps closely to this dissertation: the problem was not only to drive a simulated car, but to make alternative agent designs and their evidence inspectable by novice users.

The artefact was developed iteratively. A working TORCS runner and shared agent action contract came first, because every later comparison depends on consistent communication with the simulator. Baselines then established controllable behaviour before learning methods were introduced. Telemetry storage and GUI views were developed alongside the agents rather than added after evaluation, ensuring that the same records could support technical analysis and user-facing explanation. Finally, the Windows package made the completed application available without a development installation.

The iteration was evidence-led. When a controller failed, the response was not simply to add another algorithm. The random driver first checked the action path; the rule-based controller supplied observable stabilisation behaviour; map-aware control tested whether explicit route information improved the baseline; and the learning agents were introduced only once result recording was available. For Agent~8, early high-pace but low-completion outcomes changed the training objective from pace discovery to clean repeated completion. This reasoning makes the final system a sequence of justified experiments rather than a catalogue of unrelated features.

\subsection{Evaluation logic}

The evaluation follows a Goal-Question-Metric structure: the questions concern comparative driving behaviour and novice-facing explanation, so the metrics must capture both pace and reliability \parencite{basili1994}. The technical measures are lap completion, completed-lap time, progress, off-track behaviour, termination reason and control telemetry. The user-facing measures are functional evidence that agents can be selected, run, reviewed and compared through the GUI, plus a deployment smoke test for the packaged application.

For the learning-platform contribution, evaluation asks whether the intended novice route is actually available in the delivered artefact. The defined tasks are to launch the package, select a named agent, start a run, observe plain-language state and telemetry, retrieve a saved outcome, and compare reliability with pace. These are observable acceptance conditions, not a proxy questionnaire for user satisfaction or learning gain. They allow the dissertation to evaluate functional accessibility and explanation design honestly while retaining a participant study as the appropriate method for stronger educational claims.

No human participants or personal data were used. The principal ethical issue is the accuracy of claims. TORCS results are evidence about simulated autonomous racing, not road safety. Similarly, a release package demonstrates simpler deployment, while the absence of a novice participant study limits any claim that the application has proven learning outcomes. These boundaries were retained explicitly rather than hidden behind broad claims about artificial intelligence.

The research questions are answered by two complementary evidence streams. Product evidence, including automated tests, package smoke checks, screenshots and stored GUI runs, shows that the application integrates the required features. Performance evidence, including fixed-policy evaluation files, supplies repeated trial outcomes for learned drivers. Treating these streams separately improves validity. A dashboard screenshot can demonstrate that a comparison view exists, but it cannot prove neural-policy reliability. Conversely, an evaluation record can support a completion claim, but it cannot demonstrate that a beginner can locate and interpret it through the interface.

\subsection{Requirements and success criteria}

The requirements were grouped into technical, evidential and educational categories. Technically, the application had to launch TORCS, run several agent types, record a common set of outcomes and preserve policy files. Evidentially, it had to retain failed as well as successful runs and make their provenance inspectable. Educationally, it had to provide a graphical route through the system, use readable agent explanations and reduce deployment steps through the Windows package.

This treatment follows the principle of requirements engineering that capabilities and quality constraints should be stated in a form that can later be verified \parencite{iso29148}. It is especially useful here because a novice-friendly application can otherwise be evaluated impressionistically. By defining observable requirements, such as launching from the packaged folder, presenting an outcome after a race and retaining failures in a comparison, the project can demonstrate that the wrapper has delivered the requested route into the technical work even without claiming a completed human-subjects study.

The success criterion was therefore not \enquote{produce the fastest lap}. A successful project would provide a functioning and distributable AI racing application, show meaningful differences among control approaches, and allow a non-specialist to inspect that evidence without reading source code. This definition explains why stable engineered baselines, weak learning outcomes and failed episodes remain relevant results rather than being discarded.

The requirement table is deliberately concise because detailed implementation is assessed later. Its purpose is traceability. Each requirement has an observable acceptance condition, and each condition is linked to either a system test, stored experiment output or visible application behaviour. This establishes a chain from the brief, through design decisions, to evidence. It is stronger than declaring that the software is \enquote{user friendly} or \enquote{robust} without defining what those qualities require in this project.

